\scanpage{7}\mspage{6}
\deleted{$\mathrm{At}_{M}\,4:\ \forall X\in\unsure{M|\mathcal M},\ \unsure{\Phi|\varphi}\in\widetilde{\unsure{\mathscr F|\mathcal F|\mathfrak F}}$, posant $\Phi^{X}=\{Y\in\Phi\mid X\ll Y\}$, $\Phi^{X}$ a un plus petit \unsure{élément}\unsure{.}}

\margin{J'ai exprimé $\uline{\mathrm{At}}\,\deleted{\unsure{6}}5$ sans référence à $\mathscr F$ ou à $\widetilde{\mathscr F}$, mais directement en termes de $(M,\leq,\ll,\mathrel{\tikz[baseline=-.4ex,x=.65em,y=.65em]{\draw (0,1.1)--(1,.5)--(0,-.1) (1,1.1)--(0,.5)--(1,-.1);}})$ $\nearrow$}%
$\mathrm{At}_{M}\,4$ \quad a) \deleted{\unsure{Posant|Posons} $X\ll Y\ \overset{\unsure{\mathrm{déf}}}{\unsure{\Longleftrightarrow|\Longrightarrow}}\ \unsure{Y|X}$ \unsure{contient}} \inserted{\deleted{$\forall X,Y\in M,\ X\ll$}} \inserted{\deleted{\unsure{est|ent} \unsure{minimal|minimum}}} Si $X\ll Y$, alors il y a un \deleted{\unsure{plus petit élément}} \inserted{élément minimal} \inserted{dans} $\widetilde Y^{X}=\{Z\in\widetilde Y\mid X\ll Z\}$ \deleted{\unsure{i.e.} $X\ll Y'\leq Y\Longrightarrow Y'=Y$} \inserted{(posons $X\overset{\circ}{\ll}Y$ si c'est $Y$)}

\[
\begin{aligned}
\text{b) Alors}\quad & X\overset{\circ}{\ll}Y,\quad X\overset{\circ}{\ll}Y',\quad Y\mathrel{\tikz[baseline=-.4ex,x=.65em,y=.65em]{\draw (0,1.1)--(1,.5)--(0,-.1) (1,1.1)--(0,.5)--(1,-.1);}}Y'\ \Longrightarrow\ Y=Y'\\
\mathrm{At}_{M}\,\deleted{\unsure{5}}6\quad &\text{Si }X\mathrel{\tikz[baseline=-.4ex,x=.65em,y=.65em]{\draw (0,1.1)--(1,.5)--(0,-.1) (1,1.1)--(0,.5)--(1,-.1);}}Y,\ \text{alors }\ X\ll Y\ \Longrightarrow\ X\leq Y\\
\mathrm{At}_{M}\,\deleted{\unsure{6}}5\quad &X\overset{\circ}{\ll}Y\overset{\circ}{\ll}Z\ \Longrightarrow\ X\overset{\circ}{\ll}Z
\end{aligned}
\]
\margin{\tikz[x=.6em,y=.6em]{\draw (.95,0)..controls(-.3,-.15)and(-.2,1)..(.95,1.25); \draw[->] (.2,.95)..controls(-.18,1.15)and(0,2)..(.08,2.3)--(1.45,3);}}%

D'autre part, pour des $\mathscr F$ à figures de type fini, on peut laisser tomber l'élément de structure $\widetilde{\mathscr F}$, \inserted{au profit de $\mathrel{\tikz[baseline=-.4ex,x=.65em,y=.65em]{\draw (0,1.1)--(1,.5)--(0,-.1) (1,1.1)--(0,.5)--(1,-.1);}}$,} et remplacer $\mathrm{At}_{M}\,1$ à $\mathrm{At}_{M}\,3$ par (a) de la \deleted{$\mathrm{At}_{M}\,\unsure{1|0}\,(\mathrm{tf})$ Toute partie fermée majorée de $M$ est de t.f.} \deleted{\unsure{At}} \unsure{page|p.} 5 (Toute partie fermée majorée de $M$ est de t.f.) et la condition b sur $ {\mathrel{\tikz[baseline=-.4ex,x=.65em,y=.65em]{\draw (0,1.1)--(1,.5)--(0,-.1) (1,1.1)--(0,.5)--(1,-.1);}}}^{*}$. (page 5)
% slanted
\margin{i.e. $X\mathrel{\tikz[baseline=-.4ex,x=.65em,y=.65em]{\draw (0,1.1)--(1,.5)--(0,-.1) (1,1.1)--(0,.5)--(1,-.1);}}Y$ et $X'\leq X,\ Y'\leq Y$ $\unsure{\mathrel{\tikz[x=.5em,y=.28em]{\draw (0,.08)--(1.2,.08) (0,.42)--(1.2,.42) .. controls (1.35,.42) and (1.38,.23) .. (1.28,.2);}}|\Longrightarrow}X'\mathrel{\tikz[baseline=-.4ex,x=.65em,y=.65em]{\draw (0,1.1)--(1,.5)--(0,-.1) (1,1.1)--(0,.5)--(1,-.1);}}Y'$}%

Dans le cas général, quand on a simplement $\mathrel{\tikz[baseline=-.4ex,x=.65em,y=.65em]{\draw (0,1.1)--(1,.5)--(0,-.1) (1,1.1)--(0,.5)--(1,-.1);}}$ satisfaisant cette condition b) plus les conditions $\mathrm{At}_{M}\,4,5,6$, on trouve
\[
\begin{aligned}
\widetilde{\mathscr F}_{\mathrm{tf}}&=\text{ensemble des parties fermées }\inserted{\Phi\ \text{de t.f.}}\ \text{de }M,\leq\\
&\deleted{\text{\unsure{formées} \unsure{t.q.}}}\ \inserted{\text{telles que :}}\quad X,Y\in\Phi\Longrightarrow X\mathrel{\tikz[baseline=-.4ex,x=.65em,y=.65em]{\draw (0,1.1)--(1,.5)--(0,-.1) (1,1.1)--(0,.5)--(1,-.1);}}Y.
\end{aligned}
\]

Pour avoir $\widetilde{\mathscr F}$, il faut cependant un choix de structure en plus : se donner certains antifiltres infinis dans l'ens ordonné $\widetilde{\mathscr F}_{\mathrm{tf}}$ précédent\ldots

On peut aussi interpréter $\mathscr F$ en termes de
\[
\begin{aligned}
\mathscr F^{*}&=\{\mathrm{Mulomb}(F)\mid F\in\mathscr F\}\subset\mathrm{Fig}(M)\\
\deleted{\text{où}}\quad\mathrm{Mulomb}(F)&=\{\deleted{\unsure{X}}\ \mathrm{Omb}(\unsure{x|X})\mid X\in\widetilde F\}\\
\mathrm{Omb}\,X&=\{Y\in M\mid Y\ll X\}=M_{\ll \unsure{x|X}}.
\end{aligned}
\]
(cf Scholie GF p. 28, et axiomes, ou mieux à la \uline{page 32}), comme sous-atelier (pas néc. \unsure{spacieux|spécieux}) de $\mathrm{Fig}(M)$.
